\section{Тема 7. Арифметика}

\subsection{Постановка задачи (вариант~4)}

В массиве $X(n)$ хранятся различные вещественные числа (как
большие, так и меньшие единицы). Округлить их, оставив в каждом
по 3 значащих цифры.

\subsection{Математическое обоснование}

Для числа $x \ne 0$ найдём порядок старшего разряда:
\[
  e = \lfloor \log_{10}|x| \rfloor.
\]
Чтобы оставить 3 значащие цифры, умножим число на $10^{2-e}$,
округлим до целого и разделим обратно:
\[
  x' = \frac{\operatorname{round}\bigl(x\cdot10^{2-e}\bigr)}{10^{2-e}}.
\]
Например, для $x=0.0034567$: $e=-3$, множитель $10^{5}$, получаем
$0.00346$.

\subsection{Блок-схема алгоритма}

\begin{center}
\begin{tikzpicture}[
  node distance=10mm, start chain=going below,
  nodes={draw=black, top color=white, bottom color=lightgray!50,
         minimum width=5.6cm, align=center, font=\small},
  arrow/.style={->, >=Stealth, thick},
]
  \node[draw, rounded corners, on chain] {Start};
  \node[shape=trapezium, trapezium left angle=70,
        trapezium right angle=110, on chain, draw, join=by arrow]
       {$n$,\ $X[0..n{-}1]$};
  \node[shape=chamfered rectangle, chamfered rectangle xsep=10pt,
        on chain, draw, join=by arrow]
       (loop) {$i := 0,\ n{-}1$};
  \begin{scope}[local bounding box=scl]
    \node[shape=rectangle, on chain, draw, join=by arrow]
         {$e := \lfloor\log_{10}|X[i]|\rfloor$};
    \node[shape=rectangle, on chain, draw, join=by arrow]
         {$f := 10^{2-e}$};
    \node[shape=rectangle, on chain, draw, join=by arrow]
         {$X'[i] := \mathrm{round}(X[i]{\cdot}f)/f$};
    \node[shape=trapezium, trapezium left angle=70,
          trapezium right angle=110, on chain, draw, join=by arrow]
         {$X[i] \to X'[i]$};
    \coordinate[on chain, join=by arrow] (le);
  \end{scope}
  \draw[arrow] (scl.south) -| ($(scl.west)-(0.9,0)$) |- (loop);
  \coordinate[on chain] (sp);
  \draw[arrow] (loop) -| ($(scl.east)+(0.9,0)$) |- (sp);
  \node[draw, rounded corners, on chain, join=by arrow] {End};
\end{tikzpicture}
\end{center}

\subsection{Текст программы}

\lstinputlisting[style=Cstyle]{../src/t07.c}

\subsection{Тестовые данные}

\begin{center}
\begin{tabular}{l|l}
\toprule
Исходное число & Результат (3 значащих цифры) \\
\midrule
$123456.789$ & $123000$ \\
$0.0034567$ & $0.00346$ \\
$7.5$ & $7.5$ \\
\bottomrule
\end{tabular}
\end{center}

\textbf{Результат выполнения программы:}
\begin{lstlisting}[language={},basicstyle=\ttfamily\small,frame=single]
Введите количество элементов n: 3
Введите элементы массива:
123456.789 0.0034567 7.5
Округлённые значения (3 значащие цифры):
  123457 -> 123000
  0.0034567 -> 0.00346
  7.5 -> 7.5
\end{lstlisting}
